\begin{table}[htbp]
\centering
\caption{Penentuan solusi: alternatif, alasan, dan kelemahan yang diakui (bukti kuantitatif dipaparkan pada Bab~\ref{chap:evaluasi}).}
\label{tab:penentuan-solusi}
\small
\begin{tabular}{@{}p{2.7cm}p{6.4cm}p{4cm}@{}}
\toprule
\textbf{Keputusan} & \textbf{Alasan Dipilih} & \textbf{Kelemahan yang Diakui} \\
\midrule
\emph{Random Forest} sebagai model utama & Robust pada data \emph{sparse}, deterministik, berjalan tanpa GPU, dan interpretabel. Uji signifikansi lintas-pasien menunjukkan LSTM sedikit lebih akurat namun selisihnya \textbf{tidak bermakna klinis}, sehingga tak sepadan dengan hilangnya interpretabilitas. & LSTM sedikit lebih akurat secara statistik meski tak bermakna klinis. \\[2pt]
Gemini (API) alih-alih Ollama lokal & Tidak butuh GPU, gratis via \emph{API key}, stabil lintas platform. & Membutuhkan koneksi internet; bergantung layanan pihak ketiga. \\[2pt]
\texttt{all-MiniLM-L6-v2} (\emph{embedding}) & \emph{CPU-only}, ringan ($\sim$80~MB), tanpa \emph{server}, \emph{open-source}. & Belum \emph{fine-tuned} untuk domain diabetes Indonesia; dimensi lebih kecil dari alternatif. \\[2pt]
Streamlit & Dukungan \emph{multi-page} \& \emph{session state} native, cocok untuk \emph{slider what-if}, mudah dipelihara. & Kurang alami untuk antarmuka gaya \emph{chat}; manajemen \emph{state} manual. \\[2pt]
OhioT1DM sebagai \emph{surrogate} & Satu-satunya dataset CGM multimodal publik yang lengkap; \emph{benchmark} standar. & T1DM $\neq$ T2DM; tanpa data Indonesia. \\[2pt]
Prediksi langsung horizon tetap ($+30/+60$) & Menghindari galat kumulatif; horizon klinis bermakna; cukup untuk mengondisikan RAG. & Horizon terbatas; tidak untuk prediksi semalaman. \\[2pt]
\emph{Simplified} DT (formula, bukan ODE) & Sesuai lingkup; tak menuntut parameter pasien individual; cukup untuk \emph{what-if} kualitatif. & Bukan fisiologis penuh; tidak diklaim sebagai DT \emph{high-fidelity}. \\[2pt]
\emph{What-if} memakai formula PK, BUKAN model RF & RF bersifat korelasional; pada data latih insulin diberikan \emph{karena} glukosa tinggi, sehingga RF mempelajari arah terbalik dan memberi rekomendasi kontrafaktual yang salah arah. Formula PK mekanistik memberi arah kausal yang benar. & Formula PK menyederhanakan fisiologi; tidak ``belajar'' dari data pasien individual. \\
\bottomrule
\end{tabular}
\end{table}
